\subsection{极大平展子代数}

引理 4. —— 设 A 是 K 上一个有限次数的中心单代数，且 A 异于 K。存在 A 的一个平展（V, §6, No.~3, 第 28 页，定义 1）子代数，它异于 K。

由韦德伯恩定理（VIII, 第 120 页，定理 1），我们可以假设 \(\mathrm{A}\) 具有形式 \(\mathbf{M}_n(\mathrm{D})\)，其中 \(n\) 是一个严格正整数，\(\mathrm{D}\) 是一个以 \(\mathrm{K}\) 为中心的体。

假设 n \textgreater{} 1。元素属于 K 的对角矩阵所成的代数是 A 的一个异于 K 的平展子代数。

假设 n = 1。设 p 是 D 的特征指数。由 VIII, 第 230 页的引理 1，存在 D 的一个元素 a，使得 \(a^{p^{m}}\) 对任何自然数 m 都不属于 K。于是对充分大的 m，元素 \(x = a^{p^{m}}\) 在 K 上是可分的（V, §7, No.~7, 第 44 页，命题 13），但它不属于 K。A 的子代数 \(\mathrm{K}(x)\) 是域 K 的一个有限次数的可分扩张，因而是 K 上的一个平展子代数；它异于 K。

命题 4. —— 设 A 是 K 上一个有限次数的中心单代数。设 L 是 A 的一个子代数，L' 是 L 在 A 中的交换子。

\begin{enumerate}
\def\labelenumi{\alph{enumi})}
\item
  若 L 在 A 的半单交换子代数中是极大的，则我们有 \(L = L'\)，且 L 是 A 的一个极大交换子代数。
\item
  若 L 在 A 的平展子代数中是极大的，则我们有 \(L = L'\)，且 L 是 A 的一个极大交换子代数。
\end{enumerate}

我们知道关系式 \(L = L'\) 表示 L 是 A 的一个极大交换子代数（VIII, 第 261 页，引理 3, a)）。假设 L 是半单的、交换的且异于 \(L'\)。由 VIII, 第 259 页的定理 5，\(L'\) 是半单的，且 L 是 \(L'\) 的交换子，因而是 \(L'\) 的中心；于是，\(L'\) 不是交换的。只需证明存在 A 的一个半单交换子代数 M，它异于 L、包含 L，且当 L 是平展的时候它也是平展的。

由半单环的结构定理（VIII, 第 135 页，定理 1），存在单环 \(B_{1},\ldots,B_{r}\) 与一个同构 \(\varphi\)（从 \(L'\) 到 \(B_{1}\times\cdots\times B_{r}\) 的）。对 \(1\leqslant i\leqslant r\)，把 \(B_{i}\) 的中心记作 \(E_{i}\)；则我们有 \(\varphi(\mathrm{L})=\mathrm{E}_{1}\times\cdots\times\mathrm{E}_{r}\)。由于 \(L'\) 不是交换的，我们可以假设例如 \(B_{1}\) 不是交换的；于是我们有 \(B_{1}\neq E_{1}\)，且由引理 4，存在一个子代数 \(M_{1}\)（\(B_{1}\) 的），它是交换的、异于 \(E_{1}\) 且在 \(E_{1}\) 上是平展的。设 \(\mathrm{M} = \varphi^{-1}(\mathrm{M}_{1} \times \mathrm{E}_{2} \times \cdots \times \mathrm{E}_{r})\)；它是 A 的一个包含 L 且异于 L 的半单交换子代数。假设 L 在 K 上是平展的；我们来证明 M 是平展的。K 的扩张 \(E_{i}\) 都是可分的（V, §6, No.~4, 第 30 页，命题 3）。此外，由于 \(E_{1}\)-代数 \(M_{1}\) 与 K-代数 \(E_{1}\) 都是平展的，K-代数 \(M_{1}\) 也是平展的（V, §6, No.~5, 第 32 页，推论 2）。于是 K-代数 \(M_{1} \times E_{2} \times \cdots \times E_{r}\) 是平展的，因而 M 也是平展的。

设 A 是 K 上一个有限次数的中心单代数。A 的一个在 A 的半单交换子代数中为极大的子代数称为 A 的极大半单交换子代数。由命题 4，“极大”一词指的是交换性或半单且交换这一性质。A 的一个在 A 的平展子代数中为极大的子代数称为 A 的极大平展子代数。

推论 1. —— 设 A 是一个有限次数的中心单 K-代数。A 的每个半单（相应地，平展的）交换子代数都包含在 A 的一个半单（相应地，平展的）极大交换子代数中。

推论 2. —— 设 D 是 K 上一个有限次数的体，其中心为 K。

\begin{enumerate}
\def\labelenumi{\alph{enumi})}
\item
  D 的极大交换子域就是 D 的极大交换子代数，也是 D 的极大半单交换子代数。D 的每个交换子域 L 都包含在一个极大交换子域中。
\item
  设 L 是 D 的一个交换子域且是 K 的可分扩张；它包含在 D 的一个作为 K 的可分扩张的极大交换子域中。
\item
  设 L 是 D 的一个交换子域。则 L 是 D 的极大交换子域当且仅当我们有 \([D:K]=[L:K]^{2}\)。
\end{enumerate}

D 的子代数是一个体（V, §2, No.~2, 第 10 页，命题 1），因而是半单的。此外，D 的一个极大交换子域包含 K。于是断言 a) 与 b) 由推论 1 以及命题 3 的断言 c)（VIII, 第 262 页）得出。

命题 5. —— 设 A 是 K 上一个有限次数的中心单代数。设 B 是 A 的一个半单子代数，B' 是 B 的交换子。

\begin{enumerate}
\def\labelenumi{\alph{enumi})}
\item
  代数 B 包含 A 的一个极大半单交换子代数当且仅当 B 包含 B'。
\item
  假设 B 包含 \(B'\)，并设 g 是从 B 到 A 的一个 K-代数同态。存在 A 的一个内自同构 \(\theta\)，它在 B 上与 g 相同。
\end{enumerate}

设 L 是 A 的一个极大交换子代数；由 VIII, 第 261 页的引理 3，L 等于它在 A 中的交换子 \(L'\)。若 B 包含 L，则它的交换子 \(B'\) 包含于 \(L'\)，因而包含于 B。

反之，假设 \(B'\) 包含于 B。则 \(B'\) 是 B 的中心，且它是一个半单交换代数（VIII, 第 137 页，命题 2）。由推论 1，存在 A 的一个包含 \(B'\) 的极大半单交换子代数 L。L 的交换子是 L（VIII, 第 261 页，引理 3, a)），而 \(B'\) 的交换子等于 B（VIII, 第 259 页，定理 5）。于是关系式 \(L \supset B'\) 蕴含 \(L \subset B\)。我们证明了 a)。

我们来证明 b)。假设 B 包含 \(B'\)，并选取 A 的一个包含于 B 的极大半单交换子代数 L，它由 a) 存在。设 g 是从 B 到 A 的一个同态。由 VIII, 第 262 页的命题 3，我们有等式 \([A:K]=[L:K]^{2}\)；由 VIII, 第 263 页的推论，存在 A 的一个内自同构 \(\theta_{1}\)，它在 L 上与 g 相同。若 f 是 B 到 A 的典范单射，则同态 g 与 \(\theta_{1}\circ f\) 在 B 的中心 \(B'\) 上有相同的限制，因为 \(B'\) 包含于 L。由 VIII, 第 256 页的定理 2，存在 A 的一个内自同构 \(\theta\)，使得 \(g=\theta\circ f\)；换言之，\(\theta\) 扩张 g。

